\begin{textsample}{NEG Dim 1 – persona\_grok – Score -2.00 – t442\_grok.txt}  \label{ex:f1_neg_022}
Nestled at the intersection of Brazil 's Amazon, Cerrado, and Pantanal biomes, Serra Ricardo Franco State Park stands as a vital haven for biodiversity.
This protected area is home to unique species such as the caboclinho-do-sertão bird and the giant otter, showcasing the rich ecological tapestry of these converging ecosystems.
Established in 1997 to \textbf{safeguard} this natural treasure, the park has nonetheless suffered significant loss : 24 % of its land has been cleared for \textbf{cattle} pastures.
Farms within the park, including Paredão, have been selling cows to intermediaries like Barra Mansa, which in turn supplies major slaughterhouses operated by JBS, Minerva, and Marfrig.
These companies export beef to markets around the world, yet they have not upheld their 2009 commitments to monitor \textbf{suppliers} and exclude those linked to deforestation.
This ongoing destruction in Serra Ricardo Franco contributes to the broader surge in Amazon deforestation, heightening risks of zoonotic diseases—especially relevant in the era of COVID-19.
It 's time for stronger \textbf{enforcement} of park protections and meaningful reforms in supply chains to halt this environmental harm.

% matched lemmas: cattle, enforcement, safeguard, supplier
\end{textsample}
